
\chapter*{Synopsis}
\thispagestyle{empty}

\subsection*{Background}
Technical indicators are widely used in rule-based and algorithmic trading to
extract information about trends, momentum, and market participation from
historical price and volume data. Prior research shows that the performance of
individual indicators is often unstable across market conditions, which has led
to increased interest in combining multiple indicators within more structured
trading systems. At the same time, machine-learning-based signal filtering has
been proposed as a way to refine rule-generated trading signals while preserving
interpretability and controlling model complexity. In this thesis, these ideas
are combined in a study of a multi-indicator rule-based trading strategy applied
to the gold market through the broker symbol \texttt{XAUUSD}. 
\subsection*{Problem}
Although rule-based trading strategies are transparent and easy to interpret,
their profitability has been shown to be mixed in prior research. More complex
machine-learning approaches may improve signal selection, but they also create
a greater risk of overfitting and reduced interpretability. This creates a
relevant problem for trading strategy design: how to improve rule-based
strategies without introducing unnecessary model complexity. The thesis
addresses this problem by examining whether a simple decision-tree-based filter
can be used to screen rule-generated signals in a controlled and interpretable
way. 
\subsection*{Research Question}
The thesis investigates how the confidence of a machine-learning-based signal
evaluation model relates to the out-of-sample performance of a rule-based
trading strategy on the gold market. More specifically, the study examines whether
a decision tree can serve as a simple and interpretable model for assessing the
quality of candidate trade signals on the basis of indicator-derived features and
rule states. The purpose is not to replace the underlying rule-based strategy,
but to evaluate whether the confidence produced by a transparent machine-
learning layer can distinguish between lower-quality and higher-quality trading
signals under controlled backtesting conditions.
\subsection*{Method}
To answer the research question, the study adopts an experimental research
strategy implemented through quantitative analysis of secondary data. Historical
H1 bar data for XAUUSD are retrieved from a running MetaTrader 5 terminal via
the official MetaTrader 5 Python integration module. The analysis proceeds in
three stages. First, OHLC, volume, and spread data are transformed into
indicator-derived features corresponding to EMA, RSI, MACD, DMI, KST, and
MFI, which are used to generate rule-based buy, sell, and hold signals. Second,
signal evaluation is formulated as a supervised classification task in which each
candidate trade signal constitutes one observation, and a decision tree is trained
in Python to estimate the probability that a signal will lead to a favorable
realised outcome under the predefined exit logic. Third, the predicted
confidence scores are divided into ordered ranges, and the final strategy is
evaluated out of sample in a Python backtesting pipeline using the same market
data, transaction-cost assumptions, and exit rules across all confidence
groups.
\subsection*{Result}
The final experiment used \texttt{XAUUSD} H1 data from January 2022 to
December 2024, with a held-out out-of-sample period beginning on 28 May 2024.
The unfiltered multi-indicator baseline remained profitable out of sample, but
the strongest decision-tree confidence band outperformed that baseline on total
return, Sharpe ratio, win rate, and profit factor.
\subsection*{Discussion}
The confidence ranges did not form a perfectly increasing ordering, which means
that the decision-tree score should not be interpreted as a simple
``higher-is-better'' measure. Instead, the results suggest that the confidence
score contained useful but imperfect information for separating stronger from
weaker candidate signals in the selected \texttt{XAUUSD} H1 setting.
